\section{Introducción}\label{sec:introduccion}

La Zona Metropolitana de Monterrey (ZMM) registra sus accidentes de tránsito en
la base de Accidentes de Tránsito Terrestre en Zonas Urbanas y Suburbanas (ATUS)
del Instituto Nacional de Estadística y Geografía (INEGI) \parencite{inegi_atus}:
379,294 accidentes en 18 municipios durante 2019--2024 (§\ref{sec:datos-fuente}).
La pregunta que originó el trabajo es operativa: cuántos accidentes habrá y
dónde, de modo que la atención vial y la inversión en seguridad puedan
anticiparse a la siguiente ventana de observación.

La primera exploración mostró que la pregunta combina dos problemas de
naturaleza distinta, y la reformulación separa el «dónde» del «cuánto».

\paragraph{El «dónde» es estable.} Los accidentes se concentran en un conjunto
reducido de celdas hexagonales, y esa concentración persiste entre años y bajo
cambios de la escala de la rejilla (§\ref{sec:res-espacial}). La unidad de
análisis es la celda de la rejilla hexagonal de 500\,m de lado
(§\ref{sec:met-unidad}); se descartó el municipio como unidad porque su frontera
administrativa es arbitraria respecto del fenómeno y mezcla territorios de
tamaño muy distinto.

\paragraph{El «cuánto» depende del régimen de reporte.} El nivel de accidentes
de un municipio refleja, en parte, cómo y con qué cobertura se registran, no solo
cuántos ocurren. Los modelos de conteo muestran que el error de predicción
dominante es ese régimen municipal y no la dinámica vial
(§\ref{sec:res-conteo}); la ZMM contiene cambios de régimen abruptos dentro de
la propia serie (§\ref{sec:disc-regimen}). Con datos observacionales y sin
medida de exposición (tránsito, longitud de vialidad), además, se estima una
frecuencia y no una tasa, y los resultados se interpretan como asociaciones, no
como efectos.

Por ello la pregunta se reformuló en dos partes: predecir, por celda, el conteo
de accidentes registrados en un año futuro, y separar de ese conteo el
componente propio de la celda, neto del municipio. El segundo componente es el
efecto de celda del modelo jerárquico (§\ref{sec:res-jerarquico}); no equivale a
siniestralidad, porque mezcla riesgo real e intensidad de captura.

En la base, la coordenada se comporta como la de la
intersección reportada y no como la del punto físico del choque
(§\ref{sec:datos-calidad}); el algoritmo con que se asigna la coordenada a cada
parte no está documentado (§\ref{sec:datos-captacion}).

\paragraph{Contribuciones.} El trabajo aporta:
\begin{enumerate}
	\item una caracterización de la concentración espacial de los accidentes de la
	      ZMM sobre una partición común de 2,270 celdas, con corrección por
	      pruebas múltiples y validación de su estabilidad;
	\item un diagnóstico de que el régimen de reporte municipal domina el error de
	      los modelos de conteo, que motiva separar en el modelo un efecto
	      municipal por año y un efecto de celda neto del municipio, sin
	      cuantificar un reparto entre riesgo real y captura;
	\item un modelo jerárquico bayesiano con estructura CAR, efecto municipal
	      dinámico y efecto de celda neto del municipio, evaluado fuera de muestra
	      frente a modelos de referencia con un protocolo común
	      (§\ref{sec:res-validacion});
	\item el registro de los supuestos no contrastados y de los caminos que no
	      funcionaron (apéndice~\ref{app:supuestos},
	      §\ref{sec:disc-caminos}).
\end{enumerate}
Los resultados cuantitativos se resumen en el Resumen y en
§\ref{sec:conclusiones}.

\paragraph{Organización.} La §\ref{sec:datos} describe la fuente, su calidad y
las variables seleccionadas. La §\ref{sec:metodologia} define la unidad espacial,
el análisis exploratorio, los modelos de conteo, el modelo jerárquico y la
validación fuera de muestra. La §\ref{sec:resultados} presenta los resultados en
ese mismo orden, y la §\ref{sec:discusion} examina el régimen de reporte, el
alcance de la base, las limitaciones y los caminos descartados. La
§\ref{sec:conclusiones} cierra con las conclusiones y el trabajo futuro. El
apéndice~\ref{app:reproducibilidad} lista los comandos y \emph{notebooks} que
regeneran las cifras.
